% Part: history
% Chapter: biographies 
% Section: emmy-noether

\documentclass[../../../include/open-logic-section]{subfiles}

\begin{document}

\olfileid{his}{bio}{noe}

\olsection{एमी नोएथर}

\olphoto{noether-emmy}{एमी नोएथर}

एमी नोएथर (\textsc{नर}-टर) यांचा जन्म २३ मार्च १८८२ रोजी जर्मनीतील एरलांगेन येथे एका उच्च मध्यमवर्गीय विद्वान कुटुंबात झाला. त्यांना “आधुनिक बीजगणिताची जननी” म्हटले जाते. स्त्रियांच्या शिक्षणातील मोठ्या अडथळ्यांवर मात करत त्यांनी गणित आणि भौतिकशास्त्र या दोन्ही क्षेत्रांत मूलगामी योगदान दिले. त्या काळी जर्मनीतील मुलींनी कलांचे शिक्षण घ्यावे अशी अपेक्षा होती आणि त्यांना महाविद्यालयपूर्व शाळांमध्ये प्रवेश नव्हता. तरीही नोएथर यांनी गॉटिंगेन आणि एरलांगेन विद्यापीठांतील वर्गांत श्रोती म्हणून उपस्थिती लावली; एरलांगेन येथे त्यांचे वडील गणिताचे प्राध्यापक होते. अखेरीस १९०४ मध्ये महिला विद्यार्थिनींना प्रवेश देण्याचे एरलांगेनचे धोरण बदलल्यावर त्या तेथे विद्यार्थिनी म्हणून दाखल झाल्या. १९०७ मध्ये त्यांनी गणितात डॉक्टरेट मिळवली.

पात्रता असूनही नोएथर यांना कारकिर्दीत मोठा विरोध सहन करावा लागला. १९०८ ते १९१५ या काळात त्यांनी एरलांगेन येथे विनावेतन अध्यापन केले. याच काळात तत्कालीन अग्रगण्य गणितज्ञ डेव्हिड हिल्बर्ट यांचे त्यांच्याकडे लक्ष गेले आणि त्यांनी नोएथर यांना गॉटिंगेनला बोलावले. मात्र स्त्रियांना प्राध्यापकपदे मिळण्यास बंदी होती. त्यामुळे त्यांना हिल्बर्ट यांच्या नावाखालीच व्याख्याने देता आली, तीही विनावेतन. याच काळात त्यांनी आज 'नोएथर यांचे प्रमेय' म्हणून ओळखला जाणारा निष्कर्ष सिद्ध केला; सैद्धांतिक भौतिकशास्त्रात आजही त्याचा उपयोग होतो. अखेरीस १९१९ मध्ये त्यांना अध्यापनाचा अधिकार मिळाला. विद्यापीठातील सहकाऱ्यांचा विरोध सुरूच असताना हिल्बर्ट यांनी “महाशयांनो, विद्याशाखेची सभा स्नानगृह नव्हे,” असे उत्तर दिल्याचे सांगितले जाते.

१९२०च्या दशकाच्या उत्तरार्धात नोएथर यांनी अमूर्त बीजगणितावर लक्ष केंद्रित केले आणि त्यांच्या योगदानाने या क्षेत्रात आमूलाग्र बदल घडवला. आपल्या सिद्धतांमध्ये त्या अनेकदा तथाकथित आरोही साखळी अट वापरत. या अटीनुसार विशिष्ट संचांची काटेकोरपणे वाढत जाणारी अनंत साखळी नसते. उदाहरणार्थ, आता नोएथरियन वलये म्हणून ओळखल्या जाणाऱ्या काही बीजगणितीय रचनांत आदर्शांचा $I_1
\subsetneq I_2 \subsetneq \dots$ असा अनंत अनुक्रम नसतो. ही अट कोणत्याही अंशतः क्रमावर व्यापक करता येते (बीजगणितात आदर्शांना उपसंच-संबंधाने लावलेला क्रम हे तिचे विशेष प्रकरण आहे). तिची द्वैती अवरोही साखळी अटही विचारात घेता येते: अंशतः क्रमातील प्रत्येक काटेकोरपणे \emph{उतरता} अनुक्रम अखेरीस थांबतो. एखादा अंशतः क्रम ही अवरोही साखळी अट पूर्ण करत असेल, तर $<$ या $\Nat$ वरील क्रमाने जसे विगमन करता येते, तसे त्या क्रमानेही विगमन करता येते. अशा क्रमांना \emph{सुस्थापित} किंवा \emph{नोएथरियन} म्हणतात; त्यानुसार मिळणाऱ्या सिद्धता-तत्त्वाला \emph{नोएथरियन विगमन} म्हणतात.

नोएथर ज्यू होत्या. १९३३ मध्ये नाझी सत्तेवर आल्यावर त्यांना पदावरून काढून टाकण्यात आले. सुदैवाने त्यांना अमेरिकेतील पेनसिल्व्हेनिया राज्यातील ब्रिन मॉर येथे तात्पुरत्या पदासाठी स्थलांतर करता आले. त्या काळात त्यांनी प्रिन्स्टनमध्येही व्याख्याने दिली; मात्र ते विद्यापीठ स्त्रियांना अनुकूल नाही, असे त्यांना वाटले \citep[81]{Dick1981}. १९३५ मध्ये त्यांच्या गर्भाशयातील गाठ काढण्यासाठी शस्त्रक्रिया झाली. शस्त्रक्रियेनंतर झालेल्या संसर्गामुळे त्यांचे निधन झाले आणि ब्रिन मॉर येथे त्यांचे दफन झाले.

\begin{reading}
नोएथर यांच्या चरित्रासाठी \citet{Dick1981} पाहा. पेरिमीटर इन्स्टिट्यूट फॉर थिअरेटिकल फिजिक्सची त्यांच्या जीवनावरील आणि प्रभावावरील व्याख्याने ऑनलाइन उपलब्ध आहेत \citep{Perimeter2015}. वाचनाऐवजी ऐकायचे असल्यास, \emph{Stuff You Missed in History Class} या मालिकेतील नोएथर यांच्या जीवनावरील आणि प्रभावावरील पॉडकास्ट \citep{Frey2015} पाहा. नोएथर यांचे संकलित लेखन मूळ जर्मन भाषेत \citep{Noether1983} येथे उपलब्ध आहे.
\end{reading}

\end{document}
